% Generated by scripts/emit_appendix_f.py -- do not edit.
% Sources extracted from the live modules with inspect.getsource;
% tests/test_artifact_map.py keeps this file in sync.
\noindent\emph{Lines longer than 64 columns are wrapped with a trailing backslash for typesetting; the repository holds the byte-exact sources.}

\subsection*{The four prompt strategies (one builder per cycle)}

\noindent\texttt{scripts/llm\_benchmark.py}
{\scriptsize\begin{verbatim}
def build_zero_shot_prompt(
    context_payload: dict, attr: str, dataset_name: str, researc\
    h_ctx: str,
    **_kwargs,
) -> str:
    ctx = _single_attr_context(context_payload, attr)
    schema = _OUTPUT_SCHEMA.format(attr=attr)
    return textwrap.dedent(f"""\
    You are an AI fairness auditor. Python has pre-computed the \
    fairness metrics
    below for the "{dataset_name}" dataset. Do NOT recompute the\
     math.
    Provide qualitative analysis for the protected attribute "{a\
    ttr}".

    {AIF360_CONTEXT}

    ## Research evidence (Semantic Scholar)
    {research_ctx}

    ## Fairness context (pre-computed)
    ```json
    {ctx}
    ```

    {schema}
    """)
\end{verbatim}}

\noindent\texttt{scripts/llm\_benchmark.py}
{\scriptsize\begin{verbatim}
def build_chain_of_thought_prompt(
    context_payload: dict, attr: str, dataset_name: str, researc\
    h_ctx: str,
    **_kwargs,
) -> str:
    ctx = _single_attr_context(context_payload, attr)
    schema = _OUTPUT_SCHEMA.format(attr=attr)
    return textwrap.dedent(f"""\
    You are an AI fairness auditor. Reason step by step before w\
    riting your answer.

    Step 1 — Read each metric value for "{attr}" in "{dataset_na\
    me}".
    Step 2 — Identify which metric(s) signal the most severe dis\
    parity and why.
    Step 3 — Trace the disparity back to its structural or histo\
    rical cause.
    Step 4 — Select a mitigation algorithm from the AIF360/Fairl\
    earn list that
              specifically addresses that cause. Explain why it \
    fits.
    Step 5 — Assign severity using the DI threshold rules in the\
     context.
    Step 6 — Output your final answer as JSON only (no step-by-s\
    tep prose in output).

    {AIF360_CONTEXT}

    ## Research evidence (Semantic Scholar)
    {research_ctx}

    ## Fairness context (pre-computed)
    ```json
    {ctx}
    ```

    {schema}
    """)
\end{verbatim}}

\noindent\texttt{scripts/llm\_benchmark.py}
{\scriptsize\begin{verbatim}
def build_self_critique_prompt(
    context_payload: dict, attr: str, dataset_name: str, researc\
    h_ctx: str,
    previous_output: dict,
    **_kwargs,
) -> str:
    ctx = _single_attr_context(context_payload, attr)
    prev_json = json.dumps(previous_output, indent=2)
    schema = _OUTPUT_SCHEMA.format(attr=attr)
    return textwrap.dedent(f"""\
    You are an AI fairness auditor reviewing your own prior anal\
    ysis of "{attr}"
    in the "{dataset_name}" dataset.

    ## Your previous output
    ```json
    {prev_json}
    ```

    Self-critique checklist — assess each point, then improve:
    1. Severity: is it consistent with the DI threshold rules in\
     the context?
    2. Cause: is "why_is_wrong" specific (cites a mechanism) or \
    generic (says "bias")?
    3. Fix: does "how_to_fix" name a specific algorithm from AIF\
    360 or Fairlearn?
    4. Citations: are the research entries real and directly rel\
    evant?

    {AIF360_CONTEXT}

    ## Research evidence (Semantic Scholar)
    {research_ctx}

    ## Fairness context (pre-computed)
    ```json
    {ctx}
    ```

    {schema}
    """)
\end{verbatim}}

\noindent\texttt{scripts/llm\_benchmark.py}
{\scriptsize\begin{verbatim}
def build_constrained_prompt(
    context_payload: dict, attr: str, dataset_name: str, researc\
    h_ctx: str,
    **_kwargs,
) -> str:
    ctx = _single_attr_context(context_payload, attr)
    alg_list = "Reweighing, DisparateImpactRemover, Exponentiate\
    dGradient, ThresholdOptimizer"
    schema = _OUTPUT_SCHEMA.format(attr=attr)
    return textwrap.dedent(f"""\
    You are an AI fairness auditor operating under strict output\
     constraints.

    ## Hard constraints (violations invalidate the response)
    1. "what_is_wrong" MUST quote the exact numerical DI value f\
    rom the context.
    2. "how_to_fix" MUST name exactly one algorithm from: {alg_l\
    ist}.
    3. "severity" MUST be exactly one of: CRITICAL, HIGH, MODERA\
    TE, LOW — no other text.
    4. "supporting_research" MUST contain at least 2 entries.
    5. Do NOT speculate beyond what the provided metric values s\
    upport.

    ## Attribute: "{attr}"  |  Dataset: "{dataset_name}"

    {AIF360_CONTEXT}

    ## Research evidence (Semantic Scholar)
    {research_ctx}

    ## Fairness context (pre-computed)
    ```json
    {ctx}
    ```

    {schema}
    """)
\end{verbatim}}

\subsection*{The refusal detector (\S 7.1); the 40-character rule is the brevity operationalisation the paper reports}

\noindent\texttt{scripts/llm\_benchmark.py}
{\scriptsize\begin{verbatim}
def detect_refusal(result: dict | None, raw_text: str = "") -> b\
    ool:
    """Return True if the LLM refused or gave an unusably short \
    answer."""
    if result is None:
        return True
    lower = raw_text.lower()
    if any(phrase in lower for phrase in REFUSAL_PHRASES):
        return True
    for field in REQUIRED_RESPONSE_FIELDS:
        if not result.get(field):
            return True
    # Flag very short paragraphs (< 40 chars) as effectively emp\
    ty
    for field in ("what_is_wrong", "why_is_wrong", "how_to_fix")\
    :
        if len(str(result.get(field, ""))) < 40:
            return True
    sev = str(result.get("severity", "")).upper().strip()
    if sev and sev not in VALID_SEVERITIES:
        return True
    return False
\end{verbatim}}

\subsection*{The fabrication detector (\S 7.2): float literals absent from context, citations without a matching title 3-gram}

\noindent\texttt{scripts/llm\_benchmark.py}
{\scriptsize\begin{verbatim}
def known_metric_values(context_payload: dict) -> set[float]:
    """Collect every numeric value in the metric context (rounde\
    d for matching)."""
    values: set[float] = set()

    def walk(node: Any) -> None:
        if isinstance(node, dict):
            for v in node.values():
                walk(v)
        elif isinstance(node, list):
            for v in node:
                walk(v)
        elif isinstance(node, (int, float)) and not isinstance(n\
    ode, bool):
            values.add(float(node))

    walk(context_payload)
    return values
\end{verbatim}}

\noindent\texttt{scripts/llm\_benchmark.py}
{\scriptsize\begin{verbatim}
def detect_hallucinations(
    result: dict,
    known_titles: list[str],
    context_values: set[float] | None = None,
) -> list[str]:
    """
    Flag content not traceable to the provided context (guardrai\
    l 5:
    LLMs interpret metrics, they never assert new ones).

    Two heuristics:
    - Citations: suspicious if no 3-gram appears in any known pa\
    per title
      (case-insensitive); short single-word entries are also fla\
    gged.
    - Metric values: any decimal number in the narrative fields \
    that is not
      within tolerance of a value present in the metric context \
    (and is not a
      canonical threshold such as 0.8) is flagged as a fabricate\
    d metric.
    """
    if not result:
        return []

    flags: list[str] = []

    cited = result.get("supporting_research", [])
    known_blob = " ".join(t.lower() for t in known_titles)
    for citation in cited:
        words = citation.lower().split()
        if len(words) < 3:
            flags.append(citation)
            continue
        trigrams = [" ".join(words[i : i + 3]) for i in range(le\
    n(words) - 2)]
        if not any(tg in known_blob for tg in trigrams):
            flags.append(citation)

    if context_values is not None:
        allowed = context_values | _ALLOWED_METRIC_ANCHORS
        for field in ("what_is_wrong", "why_is_wrong", "how_to_f\
    ix"):
            for match in _METRIC_VALUE_RE.findall(str(result.get\
    (field, ""))):
                quoted = float(match)
                if not any(abs(quoted - v) <= _METRIC_TOLERANCE \
    for v in allowed):
                    flags.append(f"metric value {match} in '{fie\
    ld}' not present in provided context")

    return flags
\end{verbatim}}

\subsection*{The severity rule (\S 7.4); the CRITICAL boundary is 0.8 x 0.9 = 0.72, constructed here, with no external provenance}

\noindent\texttt{scripts/qualitative\_analysis.py}
{\scriptsize\begin{verbatim}
def classify_severity(di: float | None, threshold: float = 0.8) \
    -> str:
    if di is None:
        return "UNKNOWN"
    if di < threshold * 0.9:        # e.g. < 0.72
        return "CRITICAL"
    if di < threshold:              # e.g. < 0.80
        return "HIGH"
    if di < 0.95:
        return "MODERATE (borderline)"
    return "LOW (fair)"
\end{verbatim}}

\subsection*{The scoring rubric (\S 7.5): weights 40/20/15/15/10; note cause-alignment's flat 10.0 when the baseline lists no cause labels -- the default that scored four empty responses}

\noindent\texttt{scripts/llm\_benchmark\_common.py}
{\scriptsize\begin{verbatim}
def score_llm_output(result: dict[str, Any], baseline_payload: d\
    ict[str, Any]) -> dict[str, Any]:
    """Score a qualitative LLM output against the deterministic \
    baseline."""
    baseline_by_attr = {
        row["attribute"]: row for row in baseline_payload.get("a\
    ttributes", [])
    }
    result_by_attr = {
        row["attribute"]: row for row in result.get("qualitative\
    ", [])
    }

    completeness = _score_completeness(result, expected_attrs=li\
    st(baseline_by_attr.keys()))
    severity_agreement = _score_severity(result_by_attr, baselin\
    e_by_attr)
    cause_alignment = _score_cause_alignment(result_by_attr, bas\
    eline_by_attr)
    mitigation_specificity = _score_mitigation_specificity(resul\
    t_by_attr)
    research_grounding = _score_research_grounding(result, resul\
    t_by_attr)

    total = completeness + severity_agreement + cause_alignment \
    + mitigation_specificity + research_grounding
    return {
        "total_score": round(total, 2),
        "subscores": {
            "completeness": round(completeness, 2),
            "severity_agreement": round(severity_agreement, 2),
            "cause_alignment": round(cause_alignment, 2),
            "mitigation_specificity": round(mitigation_specifici\
    ty, 2),
            "research_grounding": round(research_grounding, 2),
        },
        "summary": build_score_summary({
            "total_score": total,
            "subscores": {
                "completeness": completeness,
                "severity_agreement": severity_agreement,
                "cause_alignment": cause_alignment,
                "mitigation_specificity": mitigation_specificity\
    ,
                "research_grounding": research_grounding,
            },
        }),
    }
\end{verbatim}}

\noindent\texttt{scripts/llm\_benchmark\_common.py}
{\scriptsize\begin{verbatim}
def _score_completeness(result: dict[str, Any], expected_attrs: \
    list[str]) -> float:
    """Score structural completeness (40 pts max).

    Supports two output formats:
    - Multi-attribute: result has 'reference_audit_spec' + 'qual\
    itative' list.
    - Single-attribute: result has 'qualitative' list with one i\
    tem; no audit spec.

    Both formats award points for the same semantic qualities — \
    identification,
    diagnosis depth, element coverage, justification depth, rese\
    arch backing,
    and attribute coverage — using format-appropriate field chec\
    ks.
    """
    score = 0.0
    qualitative = result.get("qualitative", [])
    spec = result.get("reference_audit_spec")

    if spec:
        # ── multi-attribute format ────────────────────────────\
    ────────
        if spec.get("name"):
            score += 4
        if len(spec.get("core_metrics", [])) >= 4:
            score += 8
        if len(spec.get("required_response_elements", [])) >= 4:
            score += 8
        if spec.get("why_this_spec"):
            score += 5
        if len(spec.get("supporting_research", [])) >= 2:
            score += 5
    else:
        # ── single-attribute format ───────────────────────────\
    ────────
        # Score against the first qualitative entry (there shoul\
    d be exactly one).
        item = qualitative[0] if qualitative else {}
        # 4 pts: attribute identified
        if item.get("attribute"):
            score += 4
        # 8 pts: what_is_wrong is substantive (≥100 chars ≈ "≥4 \
    metrics discussed")
        if len(str(item.get("what_is_wrong", ""))) >= 100:
            score += 8
        # 8 pts: all four narrative fields non-empty
        if all(item.get(f) for f in ("what_is_wrong", "why_is_wr\
    ong", "how_to_fix", "severity")):
            score += 8
        # 5 pts: why_is_wrong is substantive (≥100 chars ≈ justi\
    fication paragraph)
        if len(str(item.get("why_is_wrong", ""))) >= 100:
            score += 5
        # 5 pts: ≥2 supporting research entries
        if len(item.get("supporting_research", [])) >= 2:
            score += 5

    # ── attribute coverage (both formats) ─────────────────────\
    ───────
    result_attrs = {row.get("attribute") for row in qualitative}
    if set(expected_attrs).issubset(result_attrs):
        score += 10

    return min(score, 40.0)
\end{verbatim}}

\noindent\texttt{scripts/llm\_benchmark\_common.py}
{\scriptsize\begin{verbatim}
def _score_severity(result_by_attr: dict[str, dict[str, Any]], b\
    aseline_by_attr: dict[str, dict[str, Any]]) -> float:
    if not baseline_by_attr:
        return 0.0
    matches = 0
    for attr, baseline in baseline_by_attr.items():
        got = result_by_attr.get(attr, {})
        if _normalize_severity(got.get("severity", "")) == _norm\
    alize_severity(baseline.get("severity", "")):
            matches += 1
    return 20.0 * (matches / len(baseline_by_attr))
\end{verbatim}}

\noindent\texttt{scripts/llm\_benchmark\_common.py}
{\scriptsize\begin{verbatim}
def _score_cause_alignment(result_by_attr: dict[str, dict[str, A\
    ny]], baseline_by_attr: dict[str, dict[str, Any]]) -> float:
    if not baseline_by_attr:
        return 0.0
    matched = 0
    total = 0
    for attr, baseline in baseline_by_attr.items():
        why = (result_by_attr.get(attr, {}).get("why_is_wrong", \
    "") or "").lower()
        labels = baseline.get("root_cause_labels", [])
        for label in labels:
            total += 1
            plain = label.replace("_", " ")
            if plain in why:
                matched += 1
    if total == 0:
        return 10.0
    return 15.0 * (matched / total)
\end{verbatim}}

\noindent\texttt{scripts/llm\_benchmark\_common.py}
{\scriptsize\begin{verbatim}
def _score_mitigation_specificity(result_by_attr: dict[str, dict\
    [str, Any]]) -> float:
    if not result_by_attr:
        return 0.0
    hits = 0
    for row in result_by_attr.values():
        text = (row.get("how_to_fix", "") or "").lower().replace\
    ("-", "").replace("_", "")
        if any(keyword.replace("-", "").replace("_", "") in text\
     for keyword in ALGORITHM_KEYWORDS):
            hits += 1
    return 15.0 * (hits / max(len(result_by_attr), 1))
\end{verbatim}}

\noindent\texttt{scripts/llm\_benchmark\_common.py}
{\scriptsize\begin{verbatim}
def _score_research_grounding(result: dict[str, Any], result_by_\
    attr: dict[str, dict[str, Any]]) -> float:
    score = 0.0
    if len(result.get("reference_audit_spec", {}).get("supportin\
    g_research", [])) >= 2:
        score += 5
    attr_hits = 0
    for row in result_by_attr.values():
        if len(row.get("supporting_research", [])) >= 1:
            attr_hits += 1
    if result_by_attr:
        score += 10.0 * (attr_hits / len(result_by_attr))
    return score
\end{verbatim}}
